{\footnotesize\setlength{\tabcolsep}{3pt}
\begin{longtable}{>{\raggedright\arraybackslash}p{2.2cm}cccc>{\raggedright\arraybackslash}p{1.0cm}>{\raggedright\arraybackslash}p{7.2cm}}
\caption{The grid. Codes: D derived by a general rule; Dn derived only by admitting D as the head of Nom and NP; D2 derived by a rule stated once for nouns and again for determinatives (or for NP and DP); L stored on the lexeme; S stated for this construction alone; U not covered. Claims: register claims for the cell, as licensed/conditional/excluded.}\label{tab:grid}\\
\toprule
Lexeme group & I1 & I2 & I3 & I4 & Claims & Basis \\
\midrule
\endfirsthead
\toprule
Lexeme group & I1 & I2 & I3 & I4 & Claims & Basis \\
\midrule
\endhead
\bottomrule
\endfoot
\multicolumn{7}{l}{\itshape Det before a nominal (§2.1)} \\
the articles & L & L & L & L & 8/2/1 & Det use permission and target selection (a: singular count) are lexical in every account; the phrase in Det is an NP under I1, I2, and I4 and a DP under I3 (primitives table). \\
every & L & L & L & L & 4/1/2 & as the articles; the Mod use after a genitive Det is a separate lexical permission (internal Mod family). \\
no/none & L & L & L & L & 1/1/1 & the dependent form no is selected by the Det use (forms family). \\
my/mine & L & L & L & L & 1/2/1 & a genitive NP in Det, the general realization; the my/mine form selection is lexical (forms family). \\
demonstratives, such, and the wh-forms & L & L & L & L & 4/1/0 & Det use permission; agreement with the target is the demonstratives' inflection (inflection family). \\
the general quantifiers & L & L & L & L & 8/2/1 & Det use permission with target restrictions (some: plural or non-count; each: singular count). \\
the gradable four & L & L & L & L & 4/1/1 & Det use permission with target restrictions (few: plural count; much: non-count). \\
cardinals & L & L & L & L & 3/1/0 & Det use with plural targets above one; complex cardinals fill Det as phrases (§5.4). \\
compound determinatives & D & D & D & D & 1/0/0 & only as genitive NPs (someone's preferences): the general genitive-Det realization. \\
pronouns & D & D & D & D & 0/0/2 & genitive NPs (my preferences, her preferences) fill Det by the general genitive-Det realization; plain-case pronouns don't. \\
common nouns, including quantificational lot, plenty, heaps & D & D & D & D & 2/0/0 & genitive NPs (people's preferences) by the general realization; plain-case NPs (what size shoes, Sunday morning) by a restricted construction that every account states separately. \\
proper nouns & D & D & D & D & 5/0/0 & genitive NPs (Kim's preferences); Sunday morning is the restricted plain-case construction. \\
\multicolumn{7}{l}{\itshape internal Mod before a nominal (§4.3)} \\
every & L & L & L & L & 1/1/0 & Mod after a genitive Det (her every move): a lexical permission. \\
the general quantifiers & L & L & L & L & 1/0/1 & several, certain, various as modifiers (its several advantages): lexical. \\
the gradable four & L & L & L & L & 4/1/0 & the few people, her many virtues: a lexical permission for Mod use; CGEL gives few Mod after a definite determiner (p. 392). \\
cardinals & L & L & L & L & 2/0/0 & these two hundred books, the remaining three: Mod use of cardinals. \\
common nouns, including quantificational lot, plenty, heaps & D & D & D & D & 2/0/0 & nominal premodifiers (dog houses): general. \\
proper nouns & D & D & D & D & 1/0/0 & Canada Day: general nominal premodification, restricted for names. \\
adjective controls & D & D & D & D & 4/0/0 & attributive AdjPs: general. \\
\multicolumn{7}{l}{\itshape independent argument use, including the existential (§4.2)} \\
the articles & L & L & L & L & 1/1/6 & no independent argument use: the lexeme lacks the use permission in every account. \\
every & L & L & L & L & 1/0/5 & excluded by the lexeme's permissions in every account; the corpus check is tabled in the quantifier supplement. \\
no/none & D & Dn & D & D & 4/1/2 & none as the independent form (forms family); the use is licensed by the ordinary Head chain (I1, I2) or by Det–Head fusion (I3, I4). Under I2 the derivation goes through the admission of D as a head of Nom and NP. \\
my/mine & D & D & D & D & 2/0/2 & mine: ordinary Head (I1, I2) or fusion with an understood nominal (I3, I4); the agreement condition is in the agreement family. \\
demonstratives, such, and the wh-forms & D & Dn & D & D & 7/0/1 & independent this/that; which and what as interrogative heads. Under I2 the derivation goes through the admission of D as a head of Nom and NP. \\
the general quantifiers & D & Dn & D & D & 20/8/0 & some, all, both, each, several, enough independently, and the existential There are several: ordinary Head (I1, I2) or Det–Head fusion (I3, I4). Under I2 the derivation goes through the admission of D as a head of Nom and NP. \\
the gradable four & D & Dn & D & D & 11/3/0 & few and many independently, including bare few with no set under discussion (§4.2). Under I2 the derivation goes through the admission of D as a head of Nom and NP. \\
cardinals & D & Dn & D & D & 1/0/0 & independent two, ten; the almost thirty who came is in the external-determination family. Under I2 the derivation goes through the admission of D as a head of Nom and NP. \\
compound determinatives & D & Dn & D & D & 4/0/0 & someone, anyone as NPs; under I3 the DP–Nom structure of Payne 2007. Under I2 the derivation goes through the admission of D as a head of Nom and NP. \\
pronouns & D & D & D & D & 12/0/3 & the noun-headed NP without a determiner (She left). \\
common nouns, including quantificational lot, plenty, heaps & D & D & D & D & 6/1/1 & bare plurals and non-counts; singular count arguments require determination, a condition every account states. \\
proper nouns & D & D & D & D & 5/0/0 & primary naming uses (Kim left). \\
adjective controls & D & D & D & D & 8/4/0 & the rich by adjectival Mod–Head fusion, retained in all four; bare numerous, multiple, countless are attested overlaps the article leaves unanalysed (§2.3), so a U would be right for them alone. \\
\multicolumn{7}{l}{\itshape independent use with a partitive of-PP (§4.3)} \\
every & L & L & L & L & 0/0/1 & no partitive use (each of, not every of): lexical. \\
no/none & D & Dn & S & S & 0/2/1 & none of the students: the of-PP is a complement in the Nom that none heads (I1, I2); under fusion it attaches outside the DP, which the inner-phrase restriction has to state (§5.2). Under I2 the derivation goes through the admission of D as a head of Nom and NP. \\
the general quantifiers & D & Dn & S & S & 3/3/0 & some of the wine, many of them: complement in the Nom the determinative heads (I1, I2); outside the fused DP by a stated restriction (I3, I4), §5.2. Under I2 the derivation goes through the admission of D as a head of Nom and NP. \\
the gradable four & D & Dn & S & S & 2/1/1 & few of them, many of the delegates: as for the general quantifiers. Under I2 the derivation goes through the admission of D as a head of Nom and NP. \\
cardinals & D & Dn & S & S & -- & two of them: as above. Under I2 the derivation goes through the admission of D as a head of Nom and NP. \\
common nouns, including quantificational lot, plenty, heaps & D & D & D & D & 4/0/0 & a lot of the wine, a number of protesters: complement of the head noun in every account (CGEL's own analysis of a number of, pp. 351–352). \\
\multicolumn{7}{l}{\itshape external determination of an independent use (§4.3)} \\
the articles & L & L & L & L & 1/0/0 & excluded: the permissions of independent few don't transfer to the articles (§5.2). \\
every & L & L & L & L & -- & excluded, as the articles. \\
the general quantifiers & L & L & L & L & -- & excluded (the some, the all): lexical. \\
the gradable four & D & Dn & D & D & 3/2/0 & the few, the lucky few: Det plus a Nom headed by few (I1, I2); Mod–Head fusion, the adjectival configuration (I3, I4), with the exclusion of a determiner inside the fused phrase stated among the inner-phrase restrictions (§5.2). Under I2 the derivation goes through the admission of D as a head of Nom and NP. \\
cardinals & D & Dn & D & D & 0/1/0 & the two, these three, the almost thirty who came: as for few. Under I2 the derivation goes through the admission of D as a head of Nom and NP. \\
compound determinatives & L & L & L & L & -- & excluded as determinatives by the compound construction; the set-denoting a special someone is the secondary use. \\
pronouns & L & L & L & L & 1/0/0 & normally excluded; poor old me is a restricted case. \\
common nouns, including quantificational lot, plenty, heaps & D & D & D & D & 0/1/0 & the general case. \\
proper nouns & L & L & L & L & -- & restricted in primary naming uses; the Smiths and an Ophelia are the secondary use. \\
adjective controls & D & D & D & D & -- & the rich, the idle rich: adjectival Mod–Head fusion in all four. \\
\multicolumn{7}{l}{\itshape internal modifiers admitted (§4.4)} \\
the general quantifiers & L & L & L & L & 1/0/2 & degree modifiers excluded for the non-gradable determinatives (very some): lexical selection. \\
the gradable four & D & Dn & D & D & 10/0/0 & degree AdvPs select the gradable four (a lexical selection) and attach inside the Nom (I1, I2) or inside the DP (I3, I4); one permission covers the very few people and the very few who objected in every account. Under I2 the derivation goes through the admission of D as a head of Nom and NP. \\
cardinals & S & S & S & S & -- & approximatives attach internally after an external determiner (the almost thirty who came) and peripherally otherwise (almost ten): a two-site condition every account states (§4.4). \\
pronouns & L & L & L & L & 3/0/1 & poor old me: restricted AdjP premodification. \\
common nouns, including quantificational lot, plenty, heaps & D & D & D & D & 0/1/1 & productive AdjP and nominal premodification. \\
proper nouns & L & L & L & L & -- & restricted premodification (architect Norman Foster). \\
adjective controls & D & D & D & D & 3/0/1 & very tall: degree modification of adjectives. \\
\multicolumn{7}{l}{\itshape peripheral modifiers admitted (§4.4)} \\
every & D & Dn & D2 & D & 4/0/0 & almost every: NP-peripheral attachment by the general rule for peripheral modifiers, with lexical selection (I1, I2, I4); a modifier inside the DP under CGEL, so the same fact needs the DP-modifier rule beside the NP-peripheral rule that covers only you (I3), §4.4. Under I2 the derivation goes through the admission of D as a head of Nom and NP. \\
no/none & D & Dn & D2 & D & 3/0/0 & practically no: as for every. Under I2 the derivation goes through the admission of D as a head of Nom and NP. \\
the general quantifiers & D & Dn & D2 & D & 5/0/0 & nearly all, hardly any: as for every. Under I2 the derivation goes through the admission of D as a head of Nom and NP. \\
cardinals & D & Dn & D2 & D & 1/0/0 & almost ten: as for every. Under I2 the derivation goes through the admission of D as a head of Nom and NP. \\
compound determinatives & D & Dn & D2 & D & -- & hardly anyone: peripheral to the NP that anyone heads (I1, I2, I4); in DP structure under Payne 2007 (I3). Under I2 the derivation goes through the admission of D as a head of Nom and NP. \\
pronouns & D & D & D & D & 2/0/0 & only you, all we who signed up: NP-peripheral in every account. \\
common nouns, including quantificational lot, plenty, heaps & D & D & D & D & -- & only the book, almost the whole class. \\
proper nouns & D & D & D & D & -- & even Kim. \\
\multicolumn{7}{l}{\itshape relative postmodification (§2.1)} \\
demonstratives, such, and the wh-forms & D & Dn & S & S & -- & those who came, that which remains: as for few. Under I2 the derivation goes through the admission of D as a head of Nom and NP. \\
the general quantifiers & D & Dn & S & S & -- & some that I saw: as above. Under I2 the derivation goes through the admission of D as a head of Nom and NP. \\
the gradable four & D & Dn & S & S & -- & few who come ever leave: an ordinary postmodifier of the Nom that few heads (I1, I2); attaches outside the fused DP, which the inner-phrase restriction states (I3, I4). Under I2 the derivation goes through the admission of D as a head of Nom and NP. \\
cardinals & D & Dn & S & S & -- & two that I have seen: as above. Under I2 the derivation goes through the admission of D as a head of Nom and NP. \\
compound determinatives & D & Dn & S & S & -- & anyone who asks, something that you need to know: as above, and ordered after the restrictor (compound family). Under I2 the derivation goes through the admission of D as a head of Nom and NP. \\
pronouns & L & L & L & L & -- & we who have read the report: a restricted range, lexical. \\
common nouns, including quantificational lot, plenty, heaps & D & D & D & D & -- & people who came: general. \\
\multicolumn{7}{l}{\itshape the compound construction: base permissions and the post-head restrictor (§4.5)} \\
the general quantifiers & D & D & D & D & 1/0/0 & the base supplies the compound's premodifier permissions (almost anybody from almost any): inheritance from the base in every account. \\
compound determinatives & S & S & S & S & 4/0/0 & the post-head restrictor's position and non-recursion, the exclusion of external determination, and the exclusion of a pre-head adjective are conditions of the compound construction in every account (§4.5); the premodifier permissions inherited from the base (hardly any, hardly anyone) are general in all four. \\
\multicolumn{7}{l}{\itshape dependent and independent forms (no/none, my/mine) (§3.2)} \\
no/none & L & L & L & L & -- & inflectional forms of one lexeme selected by the use; the partitive requires none: lexical. \\
my/mine & L & L & L & L & -- & my/mine: as no/none. \\
pronouns & L & L & L & L & -- & she/her: case selected by function. \\
\multicolumn{7}{l}{\itshape number agreement and transparency (§2.2)} \\
my/mine & S & S & D & D & 2/0/0 & Mine is ready / Mine are ready: ordinary Head needs a constructional agreement condition (§4.6); fusion takes the number from the understood nominal (I3, I4). \\
demonstratives, such, and the wh-forms & L & L & L & L & -- & this/these agree with the target: the lexeme's inflection. \\
the general quantifiers & D & Dn & D2 & D & 1/0/0 & some of the people were: the transparency rule for quantificational heads, stated once over Nom heads (I1) or once over $\{$N, D$\}$ heads (I2); CGEL states it for the number-transparent nouns and again for the fused-head partitive (I3); under I4 the rule for nouns covers some. each of the people takes singular agreement by default, a lexical fact in every account. Under I2 the derivation goes through the admission of D as a head of Nom and NP. \\
the gradable four & D & Dn & D2 & D & -- & many of them were; a few takes plural agreement in fused-head use: as for the general quantifiers. Under I2 the derivation goes through the admission of D as a head of Nom and NP. \\
common nouns, including quantificational lot, plenty, heaps & D & D & D & D & 3/0/2 & a lot of the work was, a lot of the people were: the number-transparent nouns (CGEL §5.3.3), in every account. \\
\multicolumn{7}{l}{\itshape number and grade inflection (§2.5)} \\
demonstratives, such, and the wh-forms & L & L & L & L & 2/0/0 & number inflection of two lexemes: lexical in every account; the paradigm parallels nominal number. \\
the gradable four & L & L & L & L & 7/0/0 & grade inflection (few/fewer/fewest): lexical; shared with adjectives. \\
cardinals & L & L & L & L & -- & plural cardinals (in threes) are the secondary use (below). \\
pronouns & L & L & L & L & 0/1/0 & case and reflexive forms. \\
common nouns, including quantificational lot, plenty, heaps & D & D & D & D & -- & count nouns inflect for number: a general rule with lexical exceptions. \\
proper nouns & L & L & L & L & -- & restricted number (the Smiths) is the secondary use (below). \\
\multicolumn{7}{l}{\itshape genitive marking (§2.5)} \\
compound determinatives & D & D2 & D2 & D & 1/2/0 & everyone's, everyone else's: the head-genitive definition, 'the inflection of the head noun' (CGEL p. 479), covers a compound within Noun (I1, I4); with a separate D the definition has to be extended to D heads (I2, I3), §2.5. \\
pronouns & L & L & L & L & -- & my/mine, her: suppletive genitive forms. \\
common nouns, including quantificational lot, plenty, heaps & D & D & D & D & 1/0/0 & dog's, the King of England's: head and phrasal genitives. \\
proper nouns & D & D & D & D & 1/0/0 & Kim's, Edward's. \\
\multicolumn{7}{l}{\itshape degree modification outside the NP; comparative complements; predicative use (§5.3)} \\
the articles & D & Dn & D2 & D & 2/0/1 & the bigger the better: the as a degree modifier of comparative AdjPs, as much bigger and no better; an NP modifier (I1, I2, I4) or a DP modifier (I3). Under I2 the derivation goes through the admission of D as a head of Nom and NP. \\
no/none & D & Dn & D2 & D & 1/0/0 & no better: as the bigger. Under I2 the derivation goes through the admission of D as a head of Nom and NP. \\
the general quantifiers & D & Dn & D2 & D & 6/1/0 & enough modifying adjectives, adverbs, verbs, and PPs (good enough, quickly enough): an NP modifier under the three NP-projecting implementations, under the permission that covers a great deal smaller and lots better (I1, I2, I4); a DP modifier that needs its own permission beside the NP permission (I3), §5.3. Under I2 the derivation goes through the admission of D as a head of Nom and NP. \\
the gradable four & D & Dn & D2 & D & 4/1/0 & much bigger, a lot fewer: as for enough; the comparative complement in more than ten is licensed separately in every account (lexical), and predicative Their enemies were many is a restricted lexical fact. Under I2 the derivation goes through the admission of D as a head of Nom and NP. \\
common nouns, including quantificational lot, plenty, heaps & D & D & D & D & 14/0/1 & a great deal smaller, lots better, heaps worse, three years old: NP modifiers of adjectives in every account (CGEL pp. 549–550). \\
\multicolumn{7}{l}{\itshape predeterminer modification and the former complex determinatives (§3.1)} \\
the articles & L & L & S & L & -- & a in a few, many a, such a: the article's ordinary Det use with count-singular selection (I1, I2, I4); part of a listed unit under I3. \\
demonstratives, such, and the wh-forms & D & Dn & D & D & -- & such a mess, what a mess: degree such and exclamative what in predeterminer function (I1, I2, I4); adjectives in the same function under CGEL (I3). The difference is the category, not the machinery. Under I2 the derivation goes through the admission of D as a head of Nom and NP. \\
the general quantifiers & D & Dn & D & D & 2/0/0 & all the books, both the copies: a predeterminer modifier realized by the determinative's phrase in every account (a DP in CGEL, an NP here). Under I2 the derivation goes through the admission of D as a head of Nom and NP. \\
the gradable four & D & Dn & S & D & 0/3/0 & a few, a little, a great many: an NP in Det with its own Det, as genitive Dets already are (I1, I2, I4); many a: many in predeterminer function on an a-NP. CGEL lists a few, a little, many a, and a great many as complex determinatives with a sui generis internal structure (I3), §3.1. Under I2 the derivation goes through the admission of D as a head of Nom and NP. \\
common nouns, including quantificational lot, plenty, heaps & D & D & D & D & 1/0/0 & half a cake, and the fractions: nouns and NPs in predeterminer function in every account (CGEL p. 385 n. 25). \\
\multicolumn{7}{l}{\itshape the secondary set-denoting use (§4.5)} \\
cardinals & D & S & S & D & -- & in threes, in their hundreds: the same secondary use (I1, I4); CGEL's 'secondary use in which they inflect for number and hence belong in the noun category' is a change of category (I2, I3), §5.5. \\
compound determinatives & D & S & S & D & -- & a special someone, such an anyone, the somebodies: a secondary use within Noun, the one CGEL gives proper nouns (I1, I4); conversion of a determinative to a noun, a rule the grammar has to add (I2, I3; CGEL p. 423 n. 43), §4.5 and §5.5. \\
proper nouns & D & D & D & D & -- & an Ophelia, the Smiths: CGEL's own secondary use within Noun (p. 521). \\
\end{longtable}}
